% ============================================================================
% SECTION: RELATED WORK (CONCISE - Compressed from 1.0 to 0.65 pages)
% ============================================================================

\section{Related Work}
\label{sec:related_work}

\paragraph{Cascading Architectures.}
FrugalGPT~\cite{chen2023frugalgpt} pioneered cost-efficient LLM routing via sequential cascades: cheap models handle easy queries; expensive models process hard queries. This achieves strong cost-quality trade-offs but incurs \textbf{operational rigidity}: users often need hundreds to thousands of labeled examples to calibrate chains, train domain-specific scoring functions (typically BERT-based regressors), and manually design model orderings~\cite{chen2023frugalgpt}. Adding new models typically requires re-running the calibration dataset, creating an $O(N)$ maintenance bottleneck as model pools grow. Cascades also suffer $O(K)$ latency scaling (hard queries invoke multiple models sequentially).

\textbf{How we learn:} Cascading reliability proves effective for risk-sensitive applications. \textbf{What we address:} \ours{}'s shippable priors eliminate calibration requirements (0 user-provided examples), and online learning enables $O(1)$ model addition (30 seconds, no benchmarking). Our tiered architecture preserves cascading reliability when users enable hybrid mode, while standard mode provides $O(1)$ latency via single-shot routing.

\paragraph{Supervised Routing.}
RouteLLM~\cite{ong2024routellm} trains binary classifiers on preference data from Chatbot Arena~\cite{zheng2023judging} to route between two models. This reduces setup complexity versus FrugalGPT (no scorer training) but inherits \textbf{static limitations}: the router learns a fixed decision boundary. When new models release (e.g., DeepSeek-V3, Llama-3.3), users must collect new preference pairs at scale (typically in the thousands) and retrain the classifier---an $O(N)$ recalibration bottleneck~\cite{ong2024routellm}. RouteLLM also restricts users to binary routing (strong model vs.\ weak model), limiting cost optimization across 80+ production models.

\textbf{How we learn:} Preference-based supervision provides high-quality training signal when available. \textbf{What we address:} \ours{}'s contextual bandit decouples prompt embedding from model arms, enabling instant model addition without retraining. Users register new models via API configuration; the bandit autonomously explores their utility through online feedback. This $O(1)$ adaptation enables real-time market tracking as models evolve monthly.

\paragraph{Intent-Based Routing.}
Aurelio AI~\cite{aureliolabs2024semantic} classifies prompts into semantic intents ("math," "code," "creative"), mapping each to a designated model. This provides deterministic control but requires \textbf{manual world definition}: users must anticipate all query categories and write training utterances for each. When query distributions drift or new domains emerge, users manually update intent definitions and route mappings---an ongoing maintenance burden.

\textbf{What we address:} \ours{} automatically discovers prompt-model affinities through contextual embeddings and online learning, eliminating manual intent engineering.

\paragraph{Contextual Bandits.}
LinUCB~\cite{li2010contextual} and its variants~\cite{abbasi2011improved} provide principled online learning but can suffer substantial cold-start regret in large action spaces. Transfer learning~\cite{taylor2009transfer} and meta-learning approaches improve initialization but assume access to source tasks. \textbf{What we contribute:} Shippable priors distill offline teacher supervision into compact covariance matrices, achieving 96--99\% day-one regret reduction without requiring users to provide calibration data.

\paragraph{LLM Benchmarking.}
Holistic evaluations~\cite{liang2022holistic,srivastava2022beyond} measure model capabilities but don't predict per-prompt suitability. \ours{} uses embeddings to capture prompt difficulty, enabling contextual routing beyond aggregate benchmarks.


